\chapter{Introducción}
\label{chapter:intro}

La integración multisensorial es un proceso neuronal fundamental mediante el cual el cerebro combina información proveniente de distintas modalidades sensoriales para generar representaciones más robustas y confiables del entorno. Cuando vemos a alguien hablar, integramos automáticamente las señales visuales de los movimientos labiales con las señales auditivas de la voz, percibiendo un evento unificado más claro que cualquiera de sus componentes por separado. Este fenómeno, documentado extensamente desde los estudios pioneros en el colículo superior \citep{stein2012new}, constituye una de las capacidades más notables del sistema nervioso.

Diversos enfoques computacionales han buscado capturar los mecanismos subyacentes a esta integración. Los modelos de máxima verosimilitud proponen que el cerebro selecciona la estimación perceptual que maximiza la probabilidad de las observaciones, mientras que los modelos bayesianos incorporan conocimiento previo para obtener distribuciones posteriores sobre las variables de interés. Más recientemente, arquitecturas neuronales inspiradas en la estructura cerebral han intentado implementar estas inferencias de manera biológicamente plausible.

Entre estos desarrollos, dos trabajos resultan particularmente relevantes por representar enfoques complementarios al problema. Por un lado, el modelo de \citet{cuppini2017biologically} aborda la integración audiovisual mediante una arquitectura de tres capas jerárquicas que reproduce exitosamente fenómenos experimentales como el efecto ventrílocuo, la ilusión perceptual donde un sonido parece provenir de la ubicación de un estímulo visual simultáneo. Sin embargo, al utilizar modelos de tasa deterministas, este enfoque no captura la variabilidad y las dinámicas oscilatorias características de la actividad cortical real.

Por otro lado, el modelo de \citet{echeveste2020cortical} demuestra que redes neuronales con restricciones biológicas pueden implementar inferencia probabilística mediante muestreo estocástico, exhibiendo propiedades emergentes como oscilaciones gamma, transitorios inhibitorios y variabilidad trial-a-trial consistentes con registros electrofisiológicos. No obstante, este modelo está restringido al procesamiento de una única modalidad sensorial, sin abordar la integración multisensorial ni el problema de la inferencia causal.

Esta situación revela una brecha significativa: los modelos que capturan la arquitectura funcional de la integración multisensorial carecen de dinámicas corticales realistas, mientras que aquellos con dinámicas biológicamente plausibles están limitados al procesamiento unimodal. Además, cada modelo ha sido implementado independientemente, en distintos lenguajes y con estructuras incompatibles, dificultando su comparación sistemática y cualquier intento de integración.

La visión a largo plazo que motiva este trabajo es el desarrollo de modelos que combinen la arquitectura multimodal del enfoque de Cuppini con las dinámicas corticales del modelo de Echeveste. Un modelo así permitiría estudiar cómo fenómenos como el efecto ventrílocuo emergen de circuitos neuronales que no solo integran información de múltiples modalidades, sino que lo hacen mediante mecanismos de muestreo estocástico biológicamente plausibles. Sin embargo, alcanzar esta visión requiere primero sentar bases técnicas sólidas.

El presente trabajo representa el primer paso en esta dirección mediante la implementación del modelo de Echeveste dentro del framework Scikit-NeuroMSI \citep{paredes2025scikit}, una plataforma de código abierto en Python\footnote{\url{https://www.python.org/}} diseñada específicamente para la modelización neurocomputacional de la integración multisensorial. Esta implementación incluye la arquitectura de red con poblaciones excitatorias e inhibitorias, el modelo generativo que define el problema de inferencia, y el procedimiento de entrenamiento para muestreo.

Junto con el modelo, el trabajo desarrolla herramientas de análisis para estudiar modelos con dinámicas temporales relevantes, incluyendo visualización de la evolución de potenciales de membrana, análisis espectral para identificar oscilaciones, y comparación de estadísticas de respuesta con valores objetivo. La validación verifica que la implementación reproduce las propiedades reportadas originalmente: convergencia de momentos, variabilidad modulada por contraste, y presencia de dinámicas temporales características.

Es importante delimitar claramente el alcance de esta contribución. La fusión del modelo de Echeveste con el modelo de Cuppini, la implementación de inferencia causal multisensorial, y la validación contra datos experimentales de integración audiovisual quedan como direcciones de investigación futura. El presente trabajo se enfoca exclusivamente en establecer las bases técnicas que harán posibles estas extensiones, reconociendo que una implementación correcta y validada del modelo unimodal de Echeveste es condición necesaria para cualquier desarrollo posterior hacia el procesamiento multisensorial con dinámicas corticales realistas.

% =============================================================================
% SECCIÓN: ORGANIZACIÓN DEL TRABAJO
% =============================================================================

\section{Organización del trabajo}

Los capítulos de este trabajo se organizan de la siguiente manera:

El Capítulo~\ref{chapter:antecedentes} establece los fundamentos teóricos del trabajo, recorriendo desde el contexto histórico de la neurociencia y los mecanismos biológicos de señalización neuronal, escalando del nivel molecular a la organización cortical, hasta los principios de integración multisensorial y las técnicas de machine learning para la optimización de redes recurrentes. A continuación, el Capítulo~\ref{chapter:modelos_estudiados} detalla los dos pilares computacionales de esta investigación: el modelo de Cuppini, que define el horizonte de integración audiovisual hacia el cual se proyectan las extensiones futuras, y el modelo de Echeveste, cuya implementación de inferencia probabilística representa el núcleo del presente estudio. La infraestructura técnica se describe en el Capítulo~\ref{chapter:scikit_neuromsi}, donde se presenta Scikit-NeuroMSI. Se examina su arquitectura de software y su rol como framework de código abierto esencial para el desarrollo. Por su parte, el Capítulo~\ref{chapter:implementacion} documenta el proceso de desarrollo del modelo de Echeveste, justificando las decisiones de diseño y la creación de herramientas de análisis específicas. Finalmente, en el Capítulo~\ref{chapter:resultados} se presentan los resultados de validación, verificando que la implementación reproduce las propiedades reportadas en el trabajo original, dando paso al Capítulo~\ref{chapter:conclu}, donde se discuten las contribuciones realizadas, las limitaciones del trabajo, y las direcciones futuras de investigación hacia la integración multisensorial con dinámicas corticales realistas.